% !TeX root = ../../Book2.tex
\OptionalChapter{测不准原理}{b2:ch:20}
\BookChapterNumber{b2:ch:20}
\begin{chapterabstract}
测不准关系衡量位置与频率无法同时集中的程度。本章从积分分部推导方差乘积，完整刻画其 Gaussian 极值；再用整函数增长、有限测度的维数估计和投影核分别处理 Hardy 衰减、精确支集与近似集中。各结论按自己的假设使用，不把方差、支集和集合尾部混成一种宽度。
\end{chapterabstract}

\begin{learninggoals}\begin{itemize}
\item\label{b2:item:20-goal-1}计算位置和频率中心，核对导数中的 \(2\pi\) 常数。
\item\label{b2:item:20-goal-2}从取等条件解出 Gaussian，辨析缩放、自对偶与二次相位。
\item\label{b2:item:20-goal-3}用整函数增长证明 Hardy 临界分类。
\item\label{b2:item:20-goal-4}由投影核维数和频率平移证明有限测度支集原理。
\item\label{b2:item:20-goal-5}将近似集中的尾部误差转为集合体积下界。
\end{itemize}\end{learninggoals}

把 Gaussian 缩窄一半，其 Fourier 变换就变宽一倍。若只问能否将两个支集都限制在有限测度集合，答案却更加严格：任何非零输入都做不到。两种判断分别依赖二阶矩和精确零值，需要不同的证明。

\ProseParagraphBreak

本章使用第一编的 Plancherel、Gaussian 积分和温和分布，以及基本复分析的最大模原理、Cauchy 估计和解析唯一性。扇形增长的关键步骤在正文展开；不需要借用后章的振荡估计。Benjamin Ricaud 与 Bruno Torrésani 的正式综述\cite[Sections 3.2 and 3.3.2]{Ricaud2013Survey}可比较方差与支集集中两种观点，本地所读副本是其作者自存档，引用版本仍以2014年正式发表为准。

% !TeX root = ../../Book2.tex
\section{位置与频率}
\label{b2:sec:2001}

一个窄的脉冲需要很宽的频率带才能表示。这个判断可以指位置和频率的方差，也可以指支集或近似集中的区域；不同形式有不同假设，不能把它们当作同一条不等式。

\subsection{Fourier 对偶性}
\label{b2:subsec:200101}

本章使用单位能量归一化 \(\|f\|_2=1\)。位置测量的概率密度为 \(|f(x)|^2\)，频率测量的概率密度为 \(|\widehat f(\xi)|^2\)，二者总质量因 Plancherel 都为一。平移位置不改变频率密度，只改变变换的相位；乘 \(e^{2\pi\mathrm ib\cdot x}\) 则将频率密度平移到 \(b\)。

\ProseParagraphBreak

因此“中心在原点”是可通过对称调整的条件，“同时窄”则不是。缩放 \(f_\lambda(x)=\lambda^{d/2}f(\lambda x)\) 保持能量，把位置宽度除以 \(\lambda\)，把频率宽度乘以 \(\lambda\)。任何无尺度的测不准乘积必须在这一变换下保持不变。

\subsection{位置方差}
\label{b2:subsec:200102}

若 \(|x|f\in L^2\)，定义位置中心和方差
\[
 a=\int x|f(x)|^2\,\mathrm dx,\qquad
 \sigma_x^2\coloneqq\int|x-a|^2|f(x)|^2\,\mathrm dx.
\]
中心积分由 Cauchy–Schwarz 绝对收敛。展开平方给出 \(\int|x-c|^2|f|^2=\sigma_x^2+|c-a|^2\)，故中心确实使二阶矩最小。

\ProseParagraphBreak

位置方差不是支集半径。一个遍布全空间的 Gaussian 有有限方差；一个很小质量的远端尖峰则可以使方差很大，即使绝大部分质量集中在中心。后节的集合集中量正是为了区分这两类判断。

\subsection{频率方差}
\label{b2:subsec:200103}

若 \(|\xi|\widehat f\in L^2\)，定义频率中心 \(b=\int\xi|\widehat f|^2\) 和 \(\sigma_\xi^2=\int|\xi-b|^2|\widehat f|^2\)。这等价于 \(f\) 具有平方可积弱梯度。令 \(g(x)=e^{-2\pi\mathrm ib\cdot x}f(x)\)，则
\[
 \|\nabla g\|_2^2=4\pi^2\sigma_\xi^2.
\]
这里导数因子为 \(2\pi\)，所以测不准常数必须与本卷变换约定一致。

\ProseParagraphBreak

若 \(f\) 跳跃，弱导数可能是边界分布而不属于 \(L^2\)，频率方差便为无穷。此时下界仍可解释，但不能用一个不存在的平方梯度作积分分部。

\subsection{Heisenberg 不等式}
\label{b2:subsec:200104}

下面采用\term[heisenberg-budengshi]{Heisenberg 不等式}[Heisenberg inequality]的下述版本。

\begin{theorem}[Heisenberg]\label{b2:thm:20-heisenberg}
设 \(f\in L^2(\mathbb R^d)\)、\(\|f\|_2=1\)，位置与频率二阶矩有限。则
\[
 \sigma_x\sigma_\xi\geq\dfrac d{4\pi}.
\]
\end{theorem}
\begin{proof}
通过平移和调制，把两中心移到零。对光滑紧支集函数，
\[
 d\|f\|_2^2=-2\operatorname{Re}\int
 x f(x)\cdot\overline{\nabla f(x)}\,\mathrm dx
 \leq2\||x|f\|_2\cdot\|\nabla f\|_2.
\]
第一项由 \(\operatorname{div}(x|f|^2)\) 积分为零得到，第二项是向量 Cauchy–Schwarz。对题设一般输入，先乘趋于一的光滑截断，额外梯度项由 \(R^{-1}\|f\mathbf1_{\{|x|>R\}}\|_2\) 控制；位置加权范数也由尾部收敛。再以磨光函数逼近，梯度在 \(L^2\) 收敛，而
\[
 x(\rho_\varepsilon*f)
 =\rho_\varepsilon*(xf)+(x\rho_\varepsilon)*f
\]
使加权误差不超过趋零的常数倍 \(\|f\|_2\)。因此积分恒等式和不等式传到一般输入。代入 \(\|\nabla f\|_2=2\pi\sigma_\xi\) 即得结论。
\end{proof}
若仅考察一个坐标方向，同样的积分分部给出该方向两个标准差的乘积至少 \(1/(4\pi)\)。全维结论使用方差的和，而不是把每个方向的方差任意相乘。

% !TeX root = ../../Book2.tex
\section{极值情形}
\label{b2:sec:2002}

Heisenberg 不等式的证明经过两个估计：先取复内积的实部，再用向量 Cauchy--Schwarz。要识别极值函数，必须同时追踪这两步的等号条件。这样得到的微分方程还说明了两种不同的损失：不同方向的宽度不一致会破坏向量比例关系，二次相位则会使位置与频率产生相关。

\subsection{等号情形}
\label{b2:subsec:200201}

\Needspace{6\baselineskip}
\begin{proposition}[Heisenberg 等号情形]\label{b2:thm:20-equality}
设 \(d\geq1\)，\(f\in L^2(\mathbb R^d;\mathbb C)\)、\(\|f\|_2=1\)，且 \(|x|f\)、\(|\xi|\widehat f\) 都属于 \(L^2\)，其中 \(\widehat f\) 为负号 \(2\pi\) 归一化的 Plancherel 变换。定义
\[
 \begin{aligned}
 a&=\int_{\mathbb R^d}x|f(x)|^2\,\mathrm dx,
 \quad
 b=\int_{\mathbb R^d}\xi|\widehat f(\xi)|^2\,\mathrm d\xi,\\
 \sigma_x^2&=\int_{\mathbb R^d}|x-a|^2\cdot|f(x)|^2\,\mathrm dx,\\
 \sigma_\xi^2&=\int_{\mathbb R^d}|\xi-b|^2\cdot
       |\widehat f(\xi)|^2\,\mathrm d\xi.
 \end{aligned}
\]
则 \(\sigma_x\sigma_\xi=d/(4\pi)\)，当且仅当存在 \(a_0>0\) 和常数 \(C\in\mathbb C\)，使
\[
 f(x)=C e^{-\pi a_0|x-a|^2}e^{2\pi\mathrm i b\cdot x}
 \quad
 |C|=(2a_0)^{d/4}
\]
几乎处处成立。
\end{proposition}
\begin{proof}
先把两个中心移到零。令
\[
 g(y)=e^{-2\pi\mathrm i b\cdot(y+a)}\cdot f(y+a).
\]
平移与调制保持 \(L^2\) 范数，并将位置中心和频率中心分别移到零。有限频率二阶矩保证 \(g\in H^1(\mathbb R^d)\)，且
\[
 \||y|g\|_2=\sigma_x,
 \quad
 \|\nabla g\|_2=2\pi\sigma_\xi.
\]
在复 Hilbert 空间 \(L^2(\mathbb R^d;\mathbb C^d)\) 中，取
\[
 X(y)=y\cdot g(y),
 \quad
 Y(y)=\nabla g(y),
 \quad
 \langle X,Y\rangle
 =\int_{\mathbb R^d}X(y)\cdot\overline{Y(y)}\,\mathrm dy.
\]
这里内积对第一个变量线性。由\cref{b2:lem:20-weighted-smooth-density}，位置权重与弱梯度的平方可积性允许积分分部，得到
\[
 \operatorname{Re}\langle X,Y\rangle=-\frac d2.
\]
因此 \(X,Y\) 均非零，原不等式恰来自
\[
 \frac d2
 =-\operatorname{Re}\langle X,Y\rangle
 \leq|\langle X,Y\rangle|
 \leq\|X\|_2\cdot\|Y\|_2.
\]

若取等，第二步的等号给出 \(Y=\lambda X\)，其中 \(\lambda\) 是一个常复数。这是整个向量 Hilbert 空间中的线性相关，不是允许每个位置或每个坐标分别选择比例因子。确切地说，令
\[
 \lambda=\frac{\langle Y,X\rangle}{\|X\|_2^2},
\]
则
\[
 \|Y-\lambda X\|_2^2
 =\|Y\|_2^2-\frac{|\langle X,Y\rangle|^2}{\|X\|_2^2}=0.
\]
第一步的等号又要求 \(\langle X,Y\rangle\) 为负实数。由于
\(\langle X,Y\rangle=\overline\lambda\|X\|_2^2=-d/2\)，存在 \(c>0\)，使
\[
 \nabla g(y)=-c\,y\cdot g(y)
\]
在弱导数意义下成立。

这个弱方程已经足以确定函数，无须先假设 \(g\) 经典可微。令
\[
 u(y)=e^{c|y|^2/2}\cdot g(y).
\]
在每个有界球内，指数及其一阶导数有界，故 \(u\in H^1_{\mathrm{loc}}\)，弱乘积公式给出 \(\nabla u=0\)。固定一个球，在稍大的球内将 \(u\) 磨光；所得光滑函数在原球内梯度为零，因此为常数。磨光函数在 \(L^1\) 局部收敛到 \(u\)，这些常数也收敛，因为两个常数之差的 \(L^1\) 范数等于球的体积乘其绝对值。于是 \(u\) 在每个球上几乎处处为常数。取递增的同心球，常数在交叠处一致，得到
\[
 g(y)=C e^{-c|y|^2/2}.
\]
令 \(a_0=c/(2\pi)\)，再恢复平移与调制，就是所述形式。由
\[
 \int_{\mathbb R^d}e^{-2\pi a_0|y|^2}\,\mathrm dy
 =(2a_0)^{-d/2}
\]
和 \(\|g\|_2=1\)，可得 \(|C|=(2a_0)^{d/4}\)。

反过来，所述 Gaussian 具有有限的两个二阶矩。中心化后有
\(\nabla g=-2\pi a_0\,y g\)，故上面两步估计都取等；它的位置中心为零，频率密度也是关于零对称的 Gaussian。恢复平移与调制后，两中心正是 \(a,b\)，方差乘积仍为 \(d/(4\pi)\)。
\end{proof}

极值族允许任意位置中心、频率中心、整体相位和共同宽度，却不允许不同方向使用不同的宽度。后一个限制来自同一个向量比例常数。下一组计算会把这种限制写成协方差矩阵的关系。

\subsection{Gaussian 函数}
\label{b2:subsec:200202}

先看每个方向宽度相同的情形。对 \(t>0\)，令
\[
 f_t(x)=(2t)^{d/4}e^{-\pi t|x|^2}.
\]
一维 Gaussian 积分的乘积给出 \(\|f_t\|_2=1\)，而 Gaussian 变换公式给出
\[
 \widehat f_t(\xi)
 =(2t)^{d/4}t^{-d/2}e^{-\pi|\xi|^2/t}
 =f_{1/t}(\xi).
\]
两个密度均为偶函数，故中心为零。为计算方差，对
\[
 \int_{\mathbb R}e^{-\pi s u^2}\,\mathrm du=s^{-1/2}
\]
关于 \(s>0\) 求导，得到
\[
 \int_{\mathbb R}u^2e^{-\pi s u^2}\,\mathrm du
 =\frac1{2\pi s^{3/2}}.
\]
求导的合法性可在任意固定 \(s_0>0\) 的小邻域核验：被积函数导数的绝对值由常数倍
\(u^2e^{-\pi s_0u^2/2}\) 控制，后者可积。代入 \(s=2t\)，再对其余坐标积分，每个位置坐标的方差都是 \(1/(4\pi t)\)。利用 \(\widehat f_t=f_{1/t}\)，得到
\[
 \sigma_x^2=\frac d{4\pi t},
 \quad
 \sigma_\xi^2=\frac{dt}{4\pi},
 \quad
 \sigma_x\sigma_\xi=\frac d{4\pi}.
\]
增大 \(t\) 会缩窄位置密度，同时加宽频率密度；二者的乘积保持不变。

不同方向的宽度由实对称正定矩阵描述。若非负密度 \(p\) 满足 \(\int p=1\)、\(\int xp=0\) 及 \(\int|x|^2p(x)\,\mathrm dx<\infty\)，它的\term[xiefangcha-juzhen]{协方差矩阵}[covariance matrix]是
\(\int xx^{\mathsf T}p(x)\,\mathrm dx\)：对角元记录各坐标的方差，非对角元记录位置坐标之间的相关。这里的矩阵与稍后涉及位置和频率的相关量用途不同。

\Needspace{6\baselineskip}
\begin{proposition}[Gaussian 的协方差]\label{b2:prop:20-gaussian-covariance}
设 \(d\geq1\)，\(A\) 是 \(d\times d\) 实对称正定矩阵，令
\[
 f_A(x)=\det(2A)^{1/4}e^{-\pi x^{\mathsf T}Ax},
 \quad
 x\in\mathbb R^d.
\]
则 \(f_A\) 的 \(L^2\) 范数为一，其负号 \(2\pi\) 归一化 Fourier 变换为
\[
 \widehat f_A(\xi)
 =\det(2A)^{1/4}(\det A)^{-1/2}
       e^{-\pi\xi^{\mathsf T}A^{-1}\xi}.
\]
位置中心和频率中心均为零，协方差矩阵满足
\[
 \begin{aligned}
 \Sigma_x
 &\coloneqq\int_{\mathbb R^d}xx^{\mathsf T}|f_A(x)|^2\,\mathrm dx
 =\frac1{4\pi}A^{-1},\\
 \Sigma_\xi
 &\coloneqq\int_{\mathbb R^d}\xi\xi^{\mathsf T}
                |\widehat f_A(\xi)|^2\,\mathrm d\xi
 =\frac1{4\pi}A.
 \end{aligned}
\]
因而总方差满足
\[
 \sigma_x^2\sigma_\xi^2
 =\frac{\operatorname{tr}A\cdot\operatorname{tr}A^{-1}}{16\pi^2}
 \geq\frac{d^2}{16\pi^2},
\]
等号当且仅当 \(A=tI\)，其中 \(t>0\)。
\end{proposition}
\begin{proof}
实对称矩阵的正交对角化给出
\[
 A=O^{\mathsf T}\Lambda O,
 \quad
 \Lambda=\operatorname{diag}(\lambda_1,\ldots,\lambda_d),
 \quad
 \lambda_j>0,
 \quad
 O^{\mathsf T}O=I.
\]
在变量 \(y=Ox\) 下，Lebesgue 测度不变，二次型成为
\(\sum_j\lambda_jy_j^2\)。故
\[
 \int_{\mathbb R^d}e^{-2\pi x^{\mathsf T}Ax}\,\mathrm dx
 =\prod_{j=1}^d(2\lambda_j)^{-1/2}
 =\det(2A)^{-1/2},
\]
这正好解释归一化因子的选择。

Fourier 积分绝对收敛。在同一换元下，\(x\cdot\xi=y\cdot(O\xi)\)，所以逐坐标应用一维 Gaussian 变换公式，得到
\[
 \begin{aligned}
 \int_{\mathbb R^d}
 e^{-\pi x^{\mathsf T}Ax}e^{-2\pi\mathrm i x\cdot\xi}\,\mathrm dx
 &=\prod_{j=1}^d
   \left[\lambda_j^{-1/2}
   e^{-\pi(O\xi)_j^2/\lambda_j}\right]\\
 &=(\det A)^{-1/2}e^{-\pi\xi^{\mathsf T}A^{-1}\xi}.
 \end{aligned}
\]
乘回振幅即得变换公式；其振幅也等于
\(\det(2A^{-1})^{1/4}\)，因此 \(\widehat f_A=f_{A^{-1}}\)。

在 \(y\) 坐标中，位置密度是
\[
 \prod_{j=1}^d
 \left[(2\lambda_j)^{1/2}e^{-2\pi\lambda_jy_j^2}\right].
\]
每个因子积分为一且为偶函数，因此一阶矩及不同坐标之间的混合二阶矩为零。前面的矩计算给出第 \(j\) 个对角元为 \(1/(4\pi\lambda_j)\)。换回 \(x=O^{\mathsf T}y\)，得到
\[
 \Sigma_x
 =\frac1{4\pi}O^{\mathsf T}\Lambda^{-1}O
 =\frac1{4\pi}A^{-1}.
\]
将同一公式用于 \(\widehat f_A=f_{A^{-1}}\)，便得到
\(\Sigma_\xi=(4\pi)^{-1}A\)。

总方差等于协方差矩阵的迹，故只需估计
\[
 \left(\sum_{j=1}^d\lambda_j\right)
 \cdot
 \left(\sum_{j=1}^d\lambda_j^{-1}\right)
 \geq
 \left(\sum_{j=1}^d
       \sqrt{\lambda_j}\cdot\lambda_j^{-1/2}\right)^2
 =d^2.
\]
这是有限维 Cauchy--Schwarz；等号要求两个正向量
\((\sqrt{\lambda_j})_j\) 与 \((\lambda_j^{-1/2})_j\) 成比例，即所有 \(\lambda_j\) 相同。正交对角化随即给出 \(A=tI\)。
\end{proof}

例如在二维取 \(A=\operatorname{diag}(t,t^{-1})\)，则
\[
 \Sigma_x=\frac1{4\pi}
 \begin{pmatrix}t^{-1}&0\\0&t\end{pmatrix},
 \quad
 \Sigma_\xi=\frac1{4\pi}
 \begin{pmatrix}t&0\\0&t^{-1}\end{pmatrix},
 \quad
 \sigma_x\sigma_\xi=\frac{t+t^{-1}}{4\pi}.
\]
每一个坐标方向都达到自己的一维下界，但全维下界 \(1/(2\pi)\) 只有在 \(t=1\) 时达到。原因是求总方差时先分别相加，再取乘积：一个方向在位置上窄、另一个方向在频率上窄，并不能让两个总方差同时最小。

再看二次相位。它不改变位置密度，却会改变弱导数，因而改变频率方差。一维中可以保留复配对的虚部，得到精确记录这一增量的不等式。

\Needspace{6\baselineskip}
\begin{proposition}[含相关项的不确定性不等式]\label{b2:prop:20-correlated-uncertainty}
设 \(f\in L^2(\mathbb R;\mathbb C)\)、\(\|f\|_2=1\)，且 \(xf\)、\(\xi\widehat f\in L^2\)，其中 \(\widehat f\) 为负号 \(2\pi\) 归一化的 Plancherel 变换。令
\[
 \begin{aligned}
 a&=\int_{\mathbb R}x|f(x)|^2\,\mathrm dx,
 \quad
 b=\int_{\mathbb R}\xi|\widehat f(\xi)|^2\,\mathrm d\xi,\\
 \sigma_x^2&=\int_{\mathbb R}(x-a)^2|f(x)|^2\,\mathrm dx,\\
 \sigma_\xi^2&=\int_{\mathbb R}(\xi-b)^2
                    |\widehat f(\xi)|^2\,\mathrm d\xi.
 \end{aligned}
\]
有限频率二阶矩使弱导数 \(f'\) 属于 \(L^2\)。定义
\[
 Df=\frac1{2\pi\mathrm i}f',
 \quad
 \kappa=\operatorname{Re}\int_{\mathbb R}
       (x-a)\cdot f(x)\,\overline{(D-b)f(x)}\,\mathrm dx.
\]
则
\[
 \sigma_x^2\sigma_\xi^2
 \geq\kappa^2+\frac1{16\pi^2}.
\]
等号当且仅当存在 \(a_0>0\)、\(c\in\mathbb R\) 及常数 \(C\in\mathbb C\)，使
\[
 f(x)=C e^{-\pi(a_0+\mathrm ic)(x-a)^2}
           e^{2\pi\mathrm i bx},
 \quad
 |C|=(2a_0)^{1/4}
\]
几乎处处成立。对这些函数，
\[
 \sigma_x^2=\frac1{4\pi a_0},
 \quad
 \sigma_\xi^2=\frac{a_0^2+c^2}{4\pi a_0},
 \quad
 \kappa=-\frac{c}{4\pi a_0}.
\]
\end{proposition}
\begin{proof}
令 \(g(y)=e^{-2\pi\mathrm i b(y+a)}\cdot f(y+a)\)。它的两个中心为零，且
\[
 \|yg\|_2=\sigma_x,
 \quad
 \|Dg\|_2=\sigma_\xi,
 \quad
 \operatorname{Re}\int_{\mathbb R}yg\,\overline{Dg}\,\mathrm dy
 =\kappa.
\]
第二个等式使用 Plancherel 和 \(\widehat{Dg}(\xi)=\xi\widehat g(\xi)\)。所有配对都由 Cauchy--Schwarz 绝对可积。设
\[
 I=\int_{\mathbb R}y\cdot g(y)\cdot\overline{g'(y)}\,\mathrm dy.
\]
对 \(g\) 应用\cref{b2:lem:20-weighted-smooth-density}的一维情形，得到 \(\operatorname{Re}I=-1/2\)。注意复共轭会改变 \(D\) 中虚数因子的符号，因此
\[
 J\coloneqq\int_{\mathbb R}yg\,\overline{Dg}\,\mathrm dy
 =\frac{\mathrm i}{2\pi}I
 =\kappa-\frac{\mathrm i}{4\pi}.
\]
对 \(yg\) 和 \(Dg\) 应用复 Cauchy--Schwarz，得到
\[
 \kappa^2+\frac1{16\pi^2}
 =|J|^2
 \leq\|yg\|_2^2\cdot\|Dg\|_2^2
 =\sigma_x^2\sigma_\xi^2.
\]

由于 \(\operatorname{Im}J\neq0\)，两个向量均非零，等号当且仅当存在常复数 \(z\)，使 \(Dg=z\,yg\)。此时
\[
 \overline z\,\sigma_x^2=J
 =\kappa-\frac{\mathrm i}{4\pi},
 \quad
 \operatorname{Re}z=\frac{\kappa}{\sigma_x^2},
 \quad
 \operatorname{Im}z=\frac1{4\pi\sigma_x^2}>0.
\]
将 \(z=-c+\mathrm i a_0\) 写开，其中 \(a_0>0\)、\(c\in\mathbb R\)，弱微分方程变成
\[
 g'(y)=2\pi\mathrm i z\,y\cdot g(y).
\]
在每个有界区间中乘以光滑函数 \(e^{-\pi\mathrm i zy^2}\)，弱乘积公式说明 \(e^{-\pi\mathrm i zy^2}\cdot g(y)\) 的弱导数为零。用\cref{b2:thm:20-equality}证明中的局部磨光论证识别常数，得到
\[
 g(y)=C e^{\pi\mathrm i zy^2}
     =C e^{-\pi(a_0+\mathrm ic)y^2}.
\]
单位范数给出 \(|C|=(2a_0)^{1/4}\)，恢复中心后即得所述形式。

反过来，对这个归一化 Gaussian，
\[
 Dg=(-c+\mathrm i a_0)\,yg,
 \quad
 |g(y)|^2=(2a_0)^{1/2}e^{-2\pi a_0y^2}.
\]
其位置中心为零；由 Plancherel，
\[
 \int_{\mathbb R}\xi|\widehat g(\xi)|^2\,\mathrm d\xi
 =\int_{\mathbb R}Dg(y)\cdot\overline{g(y)}\,\mathrm dy
 =(-c+\mathrm i a_0)\int_{\mathbb R}y|g(y)|^2\,\mathrm dy=0.
\]
位置矩计算给出 \(\sigma_x^2=1/(4\pi a_0)\)。再取 \(Dg\) 的平方范数及它与 \(yg\) 的配对，得到
\[
 \sigma_\xi^2
 =|{-c+\mathrm i a_0}|^2\cdot\sigma_x^2
 =\frac{a_0^2+c^2}{4\pi a_0},
 \quad
 \kappa=-c\,\sigma_x^2=-\frac{c}{4\pi a_0}.
\]
这三个量满足所述等式。平移与调制恢复 \(a,b\)，不改变这些中心化量，证明完成。
\end{proof}

因此带二次相位的 Gaussian 都使加强式取等，而原来的 Heisenberg 下界还要求 \(\kappa=0\)，即 \(c=0\)。这里的相关量由两个 \(L^2\) 向量的复配对定义；位置密度和频率密度各在自己的变量空间中，不能把它们相乘来代替这个配对。

\subsection{Fourier 自对偶性}
\label{b2:subsec:200203}

本卷归一化下，\(e^{-\pi|x|^2}\) 的 Fourier 变换仍为自身。上面的归一化族满足 \(\widehat f_t=f_{1/t}\)，所以 \(t=1\) 恰是这个族的自对偶宽度；但每个 \(t>0\) 都是 Heisenberg 极值。自对偶要求变换前后函数相同，极值只要求位置与频率宽度的乘积最小，这两个条件并不等价。

反方向的区别也可用一个显式函数看到。只在一维，记 \(q(x)=e^{-\pi x^2}\)。它属于 Schwartz 空间且 \(\widehat q=q\)，故可逐次微分 Fourier 积分，得到
\[
 \widehat{x^nq}(\xi)
 =\left(\frac{\mathrm i}{2\pi}\right)^n q^{(n)}(\xi).
\]
对 \(q\) 求二阶和四阶导数，可直接算得
\[
 \begin{aligned}
 \widehat{x^2q}(\xi)
 &=\left(\frac1{2\pi}-\xi^2\right)q(\xi),\\
 \widehat{x^4q}(\xi)
 &=\left(\xi^4-\frac3\pi\xi^2+\frac3{4\pi^2}\right)q(\xi).
 \end{aligned}
\]
因此对
\[
 H(x)=(16\pi^2x^4-24\pi x^2+3)e^{-\pi x^2},
\]
有
\[
 \begin{aligned}
 \widehat H(\xi)
 &=\left[
 16\pi^2\left(\xi^4-\frac3\pi\xi^2+\frac3{4\pi^2}\right)
 -24\pi\left(\frac1{2\pi}-\xi^2\right)+3
 \right]q(\xi)\\
 &=(16\pi^2\xi^4-24\pi\xi^2+3)q(\xi)=H(\xi).
 \end{aligned}
\]
这里的自对偶性由两次多项式计算得到，不需要预先使用 Fourier 变换的特征函数理论。

还可以把它的方差完全算出。设
\[
 M_n=\int_{\mathbb R}x^{2n}e^{-2\pi x^2}\,\mathrm dx.
\]
Gaussian 衰减使分部积分的边界项为零，对 \(n\geq1\) 有
\[
 0=\int_{\mathbb R}
 \frac{\mathrm d}{\mathrm dx}
 \left(x^{2n-1}e^{-2\pi x^2}\right)\,\mathrm dx
 =(2n-1)M_{n-1}-4\pi M_n.
\]
结合 \(M_0=1/\sqrt2\)，这给出
\[
 M_n=\frac{(2n-1)!!}{(4\pi)^n\sqrt2},
\]
其中 \(n=0\) 时的空乘积取一。平方多项式为
\[
 (16\pi^2x^4-24\pi x^2+3)^2
 =256\pi^4x^8-768\pi^3x^6
   +672\pi^2x^4-144\pi x^2+9.
\]
逐项代入上述矩公式，得到
\[
 \|H\|_2^2=12\sqrt2,
 \quad
 \int_{\mathbb R}x^2|H(x)|^2\,\mathrm dx
 =\frac{27\sqrt2}{\pi}.
\]
于是归一化函数
\begin{equation}\label{b2:eq:20-nongaussian-selfdual}
 h(x)=\frac{(16\pi^2x^4-24\pi x^2+3)e^{-\pi x^2}}
             {\sqrt{12\sqrt2}}
\end{equation}
满足 \(\widehat h=h\) 和 \(\|h\|_2=1\)。它是偶函数，两中心为零；自对偶性又使两个方差相同，故
\[
 \sigma_x^2=\sigma_\xi^2=\frac9{4\pi},
 \quad
 \sigma_x\sigma_\xi=\frac9{4\pi}>\frac1{4\pi}.
\]
多项式在 \(x=0\) 时为正，在 \(x^2=3/(4\pi)\) 时为负，故 \(h\) 有零点并变号，不属于前述非零 Gaussian 极值族。这个例子把自对偶与极值的区别具体化：Fourier 变换保持函数本身，并不保证这个函数已经具有最小的方差乘积。

% !TeX root = ../../Book2.tex
\section{其他测不准原理}
\label{b2:sec:2003}

方差假设不能回答“两个支集能否都小”或“两个函数能否都按 Gaussian 以上的速度衰减”。本节分别给出衰减、精确支集和近似集中三种形式。

\subsection{Hardy 测不准原理}
\label{b2:subsec:200301}

\begin{lemma}[临界条带增长]\label{b2:thm:20-hardy-entire-lemma}
设整函数 \(G\) 满足 \(|G(x+\mathrm iy)|\leq C e^{Ax^2}\)，且实轴上 \(|G(x)|\leq C\)，其中 \(A>0\)。则 \(G\) 为常数。
\end{lemma}
\begin{proof}
先记录扇形最大值方法。若扇形开角 \(\theta_0<\pi/2\)，全纯函数在两条边界射线有界，且增长至多 \(e^{CR^2}\)，选 \(2<\rho<\pi/\theta_0\)，乘
\(\exp(-\delta e^{-\mathrm i\rho\theta_0/2}z^\rho)\)。
在扇形内所减指数的实部至少为正的常数倍 \(\delta|z|^\rho\)，外圆弧的增长因而消失。对截断扇形用最大模原理，令外半径趋无穷，再令 \(\delta\downarrow0\)，得到内部同样有界。分数幂在该扇形选择连续辐角；原点可用小圆弧逼近处理。

现在在第一象限取 \(G_\varepsilon(z)=G(z)e^{\mathrm i\varepsilon z^2}\)。其模不超过 \(C e^{Ax^2-2\varepsilon xy}\)。选择 \(\theta_0=\arctan(A/(2\varepsilon))<\pi/2\)，在辐角至少 \(\theta_0\) 的区域已有界；在实轴和射线 \(\theta_0\) 上也有界。上述扇形论证补齐剩下区域。因此 \(|G(z)|\leq C e^{2\varepsilon xy}\)。令 \(\varepsilon\downarrow0\)，第一象限的界为 \(C\)。其他象限选择相应正负二次相位重复，得到全平面有界。Liouville 定理使它为常数；该定理由 Cauchy 导数估计 \(|G'(z)|\leq C/R\) 令 \(R\to\infty\) 得到。
\end{proof}
下面采用\term[hardy-buquedingxing]{Hardy 不确定性定理}[Hardy uncertainty theorem]的下述版本。

\begin{theorem}[Hardy]\label{b2:thm:20-hardy}
设 \(f\) 可测，\(|f(x)|\leq C e^{-\pi a|x|^2}\)、\(|\widehat f(\xi)|\leq C e^{-\pi b|\xi|^2}\)，其中 \(a,b>0\)。若 \(ab>1\)，则 \(f=0\) 几乎处处；若 \(ab=1\)，则 \(f(x)=c e^{-\pi a|x|^2}\)；若 \(ab<1\)，存在无穷多线性无关的非零输入满足这些界。
\end{theorem}
\begin{proof}
先取一维。Gaussian 支配使 \(F(z)=\int f(x)e^{-2\pi\mathrm ixz}\,\mathrm dx\) 为整函数，配方给出 \(|F(x+\mathrm iy)|\leq C'e^{\pi y^2/a}\)。当 \(b=1/a\) 时，\(G(z)=e^{\pi z^2/a}F(z)\) 满足上一引理的增长与实轴界，故 \(F(z)=c'e^{-\pi z^2/a}\)。反演得到题设 Gaussian。若 \(b>1/a\)，先用较弱的临界界推导同一形式，再比较严格更快的频率衰减，迫使常数为零。

高维中 \(F\) 在每个复变量全纯，其增长为 \(C'e^{\pi|\operatorname{Im}z|^2/a}\)。固定其他实变量，逐个使用一维引理，使 \(F(\xi)e^{\pi|\xi|^2/a}\) 对每个坐标为常数，得到同样结论。若 \(ab<1\)，选择 \(a<c<1/b\)，任何多项式乘 \(e^{-\pi c|x|^2}\) 都满足位置界，其 Fourier 变换为多项式乘 \(e^{-\pi|\xi|^2/c}\)，也满足严格留有余量的频率界。这些输入线性无关。
\end{proof}
这里使用的是点态 Gaussian 上界，不是仅有二阶矩。临界取等分类与方差极值相呼应，但两个定理不能用同一组假设相互替代。

\subsection{支集型原理}
\label{b2:subsec:200302}

若 \(f\) 的位置支集紧，Fourier 变换可延拓为整函数。若频率支集也紧，变换在某个实开集为零，解析唯一性给出 \(f=0\)。这个经典论证要求紧支集；有限测度集合可能无界，不能据此声称已经证明有限测度版本。
下面采用\term[benedicks-dingli]{Benedicks 定理}[Benedicks theorem]的下述版本。

\begin{theorem}[Benedicks]\label{b2:thm:20-benedicks}
若 \(f\in L^2(\mathbb R^d)\)，且 \(f\) 与 \(\widehat f\) 的本性支集都有有限 Lebesgue 测度，则 \(f=0\)。
\end{theorem}
\begin{proof}
记两个支集为 \(E,F\)。先证明同时支撑于 \(E\) 和频率集合 \(U\) 的 \(L^2\) 空间维数不超过 \(|E|\cdot|U|\)。算子 \(A=\mathbf1_E\mathcal F^{-1}\mathbf1_U\mathcal F\) 的 Hilbert–Schmidt 范数平方为 \(|E|\cdot|U|\)，因为从频率到位置的核为 \(\mathbf1_E(x)e^{2\pi\mathrm ix\cdot\xi}\mathbf1_U(\xi)\)。若 \(g_1,\ldots,g_N\) 是上述空间的标准正交族，\(Ag_j=g_j\)，于是 \(N\leq\|A\|_{\mathrm{HS}}^2\)。

假设 \(f\ne0\)。由示性函数的 \(L^1\) 平移连续性，可以选取互异实数 \(t_j\to0\)，使 \(|(F+t_je_1)\setminus F|<2^{-j}\)。令 \(U=F\cup\bigcup_j(F+t_je_1)\)，则 \(|U|\leq|F|+1\)。所有 \(e^{2\pi\mathrm it_jx_1}f(x)\) 都支撑于 \(E\)，Fourier 支集包含于 \(U\)。它们线性无关：若有限线性组合为零，则指数多项式在 \(\{f\ne0\}\) 上为零；Fubini 说明其一维零集有正测度。一维解析唯一性使指数多项式恒零，再用互异频率的 Vandermonde 导数矩阵得到全部系数为零。因此上述空间无限维，与有限 Hilbert–Schmidt 界矛盾。
\end{proof}
这个证明先以积分核限制空间维数，再用很小的频率平移制造任意多个独立输入。集合的有限测度只在核平方积分和平移连续性中使用，不要求其边界规则。

\subsection{定性形式}
\label{b2:subsec:200303}

对有限测度 \(E,F\)，定义位置投影 \(P_Ef=\mathbf1_Ef\) 与频率投影 \(Q_Ff=\mathcal F^{-1}(\mathbf1_F\widehat f)\)。前段说明 \(P_EQ_F\) 是紧算子。实际上其算子范数 \(\rho\) 严格小于一：若等于一，紧算子会取得范数，两个正交投影的等号条件给出一个非零函数同时具有上述两个支集，与 Benedicks 矛盾。
\begin{corollary}\label{b2:thm:20-annihilating-pair}
存在依赖 \(E,F\) 的有限常数，使
\[
 \|f\|_2\leq C(E,F)\left(
 \|f\mathbf1_{E^c}\|_2+\|\widehat f\mathbf1_{F^c}\|_2\right).
\]
\end{corollary}
\begin{proof}
有 \(\|f-P_EQ_Ff\|_2\leq\|f-P_Ef\|_2+\|P_E(f-Q_Ff)\|_2\)，后一项不超过频率尾部。另一方面 \(\|f-P_EQ_Ff\|_2\geq(1-\rho)\|f\|_2\)，取 \(C=(1-\rho)^{-1}\)。
\end{proof}
这是依赖具体集合的定量结论，常数没有被估计成只依赖体积。Benjamin Ricaud 与 Bruno Torrésani 的《A Survey of Uncertainty Principles and Some Signal Processing Applications》\cite[Section 3.3.2]{Ricaud2013Survey}讨论支集集中与这类不等式；本章的维数论证独立给出有限测度版本。

\subsection{定量形式}
\label{b2:subsec:200304}

下面采用\term[donoho-stark-budengshi]{Donoho--Stark 不等式}[Donoho--Stark inequality]的下述版本。

\begin{theorem}[Donoho--Stark 型估计]\label{b2:thm:20-donoho-stark}
设 \(\|f\|_2=1\)，有限测度集合 \(E,F\) 满足
\[
 \|f\mathbf1_{E^c}\|_2\leq\varepsilon_E,\qquad
 \|\widehat f\mathbf1_{F^c}\|_2\leq\varepsilon_F.
\]
则 \(|E|\cdot|F|\geq(1-\varepsilon_E-\varepsilon_F)_+^2\)。
\end{theorem}
\begin{proof}
三角不等式给出 \(\|P_EQ_Ff\|_2\geq1-\varepsilon_E-\varepsilon_F\)。核平方积分又给出 \(\|P_EQ_F\|_{2\to2}\leq\sqrt{|E|\cdot|F|}\)，合并得到结论。
\end{proof}
它只记录面积乘积，因此比依具体集合的严格谱界粗一些，却可以直接用于近似频带和观测窗口。若两个误差之和超过一，下界退化为零；这不是失败，而是输入尚未同时高度集中。

\ProseParagraphBreak

熵、分数阶矩和观测集合上的传播估计提供其他定量形式。使用时必须重新列出空间、归一化与指数，不能把 Heisenberg 的常数移植过去。Gerald B. Folland 与 Alladi Sitaram 的《The Uncertainty Principle: A Mathematical Survey》\cite{Folland1997Uncertainty}提供较广的正式综述；此处不把综述中未展开的全部形式列为已经证明的结论。


\begin{chaptersummary}
方差下界来自位置乘法与导数之间的积分恒等式，取等把函数限制为无二次相位的各向同性 Gaussian。Gaussian 衰减条件则经整函数增长产生更刚性的分类。支集问题用有限测度核限制维数，近似集中用投影范数控制误差。三种方法记录的集中程度不同，各自的端点与常数只能在相应假设下解释。
\end{chaptersummary}

\BookExercises

% 自拟：围绕本章已经建立的工具，改变条件或组合方法。
\begin{optionalexercise}\label{b2:ex:20-new-1}
设 \(f(x)=Ce^{-\pi(a+\mathrm ic)x^2}\)，\(a>0,c\in\mathbb R\)，归一化能量。计算两个方差和其乘积，解释二次相位的作用。
\end{optionalexercise}
\begin{solution}
位置密度正比于 \(e^{-2\pi ax^2}\)，所以 \(\sigma_x^2=1/(4\pi a)\)。求导 \(f'=-2\pi(a+\mathrm ic)xf\)，Plancherel 得 \(\sigma_\xi^2=(a^2+c^2)/(4\pi a)\)。乘积为 \(\sqrt{a^2+c^2}/(4\pi a)\)，恰在 \(c=0\) 时取下界。二次相位不改变位置密度，却增加频率方差。
\end{solution}

% 自拟：围绕本章已经建立的工具，改变条件或组合方法。
\begin{optionalexercise}\label{b2:ex:20-new-2}
对正定实对称 \(A\)，证明 \(\operatorname{tr}A\cdot\operatorname{tr}A^{-1}\geq d^2\)，并解释各向异性 Gaussian 的取等条件。
\end{optionalexercise}
\begin{solution}
设正特征值为 \(\lambda_j\)，Cauchy–Schwarz 给出 \((\sum\lambda_j)(\sum\lambda_j^{-1})\geq(\sum1)^2\)。等号要求 \(\sqrt{\lambda_j}\) 与 \(\lambda_j^{-1/2}\) 成比例，故所有特征值相同。代入正文协方差公式得测不准下界，取等恰为标量矩阵。
\end{solution}

% 自拟：围绕本章已经建立的工具，改变条件或组合方法。
\begin{optionalexercise}\label{b2:ex:20-new-3}
设 \(\|f\|_2=1\) 且两方差有限。计算平移、调制和能量保持缩放对中心与方差的影响。
\end{optionalexercise}
\begin{solution}
平移 \(f(x-a_0)\) 将位置中心加 \(a_0\)，频率中心不变，两个方差不变。调制将频率中心加 \(b_0\)，位置中心不变，方差仍不变。缩放 \(\lambda^{d/2}f(\lambda x)\) 将位置中心与标准差除以 \(\lambda\)，将频率中心与标准差乘以 \(\lambda\)。逐项换元得到这些式子，乘积保持。
\end{solution}

% 自拟：围绕本章已经建立的工具，改变条件或组合方法。
\begin{optionalexercise}\label{b2:ex:20-new-4}
说明区间示性函数的频率方差为何无穷，不能把其有限支集长度直接代入方差证明中的有限梯度公式。
\end{optionalexercise}
\begin{solution}
其变换为 \(e^{-\pi\mathrm i\xi}\sin(\pi\xi)/(\pi\xi)\)。乘 \(\xi^2\) 后平方模为 \(\sin^2(\pi\xi)/\pi^2\)，全积分发散。弱导数为 \(\delta_0-\delta_1\)，不是 \(L^2\) 函数。有限位置支集与有限频率二阶矩是独立要求。
\end{solution}

% 自拟：围绕本章已经建立的工具，改变条件或组合方法。
\begin{optionalexercise}\label{b2:ex:20-new-5}
设 \(ab<1\)。构造两种满足 Hardy 上界而不是彼此常数倍的输入，核对位置和频率两侧都留有衰减余量。
\end{optionalexercise}
\begin{solution}
取 \(a<c<1/b\)，用 \(f_0=e^{-\pi c|x|^2}\) 与 \(f_1=x_1e^{-\pi c|x|^2}\)。位置上的多项式被 \(e^{-\pi(c-a)|x|^2}\) 吸收；频率上 \(f_1\) 的变换是线性多项式乘 \(e^{-\pi|\xi|^2/c}\)，由 \(1/c>b\) 吸收。一个偶、一个在首坐标奇，故非比例。
\end{solution}

% 自拟：围绕本章已经建立的工具，改变条件或组合方法。
\begin{optionalexercise}\label{b2:ex:20-new-6}
在 \(\mathbb C^N\) 上使用单位归一化离散 Fourier 变换。证明任意非零向量的两个支集基数乘积至少为 \(N\)，并在 \(N=mn\) 时给出取等的非平凡例子。
\end{optionalexercise}
\begin{solution}
归一化输入能量为一，若位置支集基数为 \(s\)，则每个变换分量的模不超过 \(\sqrt{s/N}\)。若频率支集基数为 \(t\)，平方求和给出 \(1\leq st/N\)。取位置支集为循环群中的子群 \(\{0,n,2n,\ldots,(m-1)n\}\)，函数在其上常数。几何级数说明变换只在其零化子 \(\{0,m,2m,\ldots,(n-1)m\}\) 非零，因此支集乘积为 \(mn=N\)。
\end{solution}

% 自拟：围绕本章已经建立的工具，改变条件或组合方法。
\begin{optionalexercise}\label{b2:ex:20-new-7}
设位置尾部和频率尾部分别不超过 \(1/10\) 与 \(1/5\)。用集合集中估计给出体积乘积下界，并说明该下界没有禁止大体积集合中的近似集中。
\end{optionalexercise}
\begin{solution}
误差之和为 \(3/10\)，所以 \(|E|\cdot|F|\geq(7/10)^2=49/100\)。该估计只排除更小的体积乘积，允许更大的区域承载大部分能量。精确有限测度双支集则由 Benedicks 完全排除，这是不同的零误差刚性结论。
\end{solution}

% 自拟：围绕本章已经建立的工具，改变条件或组合方法。
\begin{optionalexercise}\label{b2:ex:20-new-8}
设 \(E,F\) 有限测度，\(\rho=\|P_EQ_F\|<1\)。证明若一列单位能量输入的位置尾部与频率尾部同时趋零，就与固定集合的谱界矛盾。
\end{optionalexercise}
\begin{solution}
定量湮灭对估计给出 \(1\leq(1-\rho)^{-1}(\varepsilon_{E,n}+\varepsilon_{F,n})\)，右侧趋零，矛盾。即使 \(|E|\cdot|F|>1\)，粗面积估计可能没有矛盾，固定集合的严格谱界仍排除这种序列。这解释了常数依赖集合形状的意义。
\end{solution}
